\documentclass[11pt]{article}
\usepackage[margin=1in]{geometry}
\usepackage{graphicx}
\usepackage{booktabs}
\usepackage{longtable}
\usepackage{hyperref}
\usepackage{float}
\usepackage{caption}
\usepackage[table]{xcolor}

\title{Synthetic ERCOT Grid --- Substation Design Summary}
\author{Texas A\&M University --- Overbye Research Group}
\date{\today}

\begin{document}
\maketitle

%% ============================================================
\section{System Overview}
%% ============================================================

\begin{table}[H]
\centering
\caption{System Overview}
\begin{tabular}{lr}
\toprule
\textbf{Metric} & \textbf{Value} \\
\midrule
Total Substations          & 4,519 \\
Substation with load/gen   & 4,399 \\
Large-Load Substations ($>$75 MW) & $\sim$250 (tunable) \\
Total Buses                & 9,434 \\
Avg Buses per Substation   & 2.09 \\
Total System Load          & 145,572 MW \\
Total Generation Capacity  & $\sim$268,425 MW \\
Weather Zones              & 8 \\
Areas                      & 8 \\
Zones                      & ? \\
Counties                   & 200 \\
\bottomrule
\end{tabular}
\end{table}

%% ============================================================
\section{Data Sources}
%% ============================================================

\begin{table}[H]
\centering
\caption{Data Sources}
\small
\begin{tabular}{p{3.2cm}p{4.5cm}p{6.5cm}}
\toprule
\textbf{Data} & \textbf{Source} & \textbf{Description} \\
\midrule
Substation locations & U.S. Census Bureau ACS 5-Year (2022) & Census tract centroids used as load fragment locations within ERCOT boundary \\
Operable generation & EIA Form 860 (2024) & Plant coordinates, generator capacity, fuel type, and operational status \\
Planned generation & ERCOT GIS Interconnection Queue (Jan 2026) & IA-signed approved projects from ERCOT GIS Report \\
Weather zones & ERCOT Weather Zone Shapefile & 8 ERCOT climate zones used for area assignment \\
County boundaries & U.S. Census Bureau TIGER/Line Shapefiles & Texas county polygons used for zone assignment \\
Reference case & 2026 ERCOT 70k-bus PowerWorld model (ver 30, 2025-11-14) & Bus kV distribution and load patterns used as statistical targets (ONLY FOR VALIDATION) \\
Large load spatial pattern & ERCOT 70k-bus Reference Case & Loads $\geq$ 75 MW extracted and fit to Thomas Cluster Process \\
Data center locations & IM3 Open Source Data Center Atlas & TX data center coordinates for DataCenter labeling (50 km proximity) \\
765 kV corridors & ERCOT Permian Basin Transmission Plan & Permian Basin and Eastern corridor waypoints (offset 20--50 km) \\
System load target & 2031 ERCOT CDR Forecast & 145 GW total system load including $\sim$52 GW large loads \\
\bottomrule
\end{tabular}
\end{table}

%% ============================================================
\section{Weather Zones and Counties}
%% ============================================================

\begin{figure}[H]
\centering
\includegraphics[width=0.85\textwidth]{figures/weather_zones_counties.png}
\caption{ERCOT Weather Zones and Texas Counties}
\end{figure}

%% ============================================================
\section{Substation Composition}
%% ============================================================

\begin{table}[H]
\centering
\caption{Substation Composition}
\begin{tabular}{lrr}
\toprule
\textbf{Type} & \textbf{Count} & \textbf{\% of Core} \\
\midrule
Load only        & 3,568 & 81.1\% \\
Gen only         & 240   & 5.5\%  \\
Both load + gen  & 591   & 13.4\% \\
EHV ($>$200 kV)  & 1,590 & 36.1\% \\
\bottomrule
\end{tabular}
\end{table}

%% ============================================================
\section{Voltage Level Distribution}
%% ============================================================

\begin{table}[H]
\centering
\caption{Voltage Level Distribution by Buses}
\begin{tabular}{rrr}
\toprule
\textbf{kV} & \textbf{Synthetic} & \textbf{Reference case} \\
\midrule
69  & 1,962 (21.1\%) & 10,241 (15.3\%) \\
138 & 5,798 (62.3\%) & 46,191 (69.2\%) \\
345 & 1,542 (16.6\%) & 10,287 (15.4\%) \\
765 & 12 (0.1\%)     & 0 (0.0\%) \\
\midrule
\textbf{Total} & \textbf{9,314} & \textbf{66,719} \\
\bottomrule
\end{tabular}
\end{table}

\noindent\textit{Reference case: 2026 ERCOT 70k-bus PowerWorld model (ver 30, 2025-11-14)}

\begin{figure}[H]
\centering
\includegraphics[width=0.85\textwidth]{figures/765kV_design_overview_panel4.png}
\caption{Voltage Level Distribution}
\end{figure}

%% ============================================================
\section{Load Distribution by Weather Zone}
%% ============================================================

\begin{table}[H]
\centering
\caption{Load Distribution by Weather Zone}
\begin{tabular}{lrr}
\toprule
\textbf{Zone} & \textbf{Share} & \textbf{MW} \\
\midrule
North Central  & 26.6\% & 38,722 \\
Coast          & 23.5\% & 34,210 \\
South Central  & 14.1\% & 20,525 \\
Far West       & 12.4\% & 18,051 \\
West           & 9.4\%  & 13,684 \\
South          & 7.4\%  & 10,772 \\
North          & 3.7\%  & 5,386 \\
East           & 3.1\%  & 4,513 \\
\bottomrule
\end{tabular}
\end{table}

\noindent\textit{Reference: Load fragments derived from U.S. Census ACS 2022 tract data, scaled to 2031 ERCOT CDR target ($\sim$145 GW)}

\begin{figure}[H]
\centering
\includegraphics[width=0.85\textwidth]{figures/765kV_design_overview_panel1.png}
\caption{Load Distribution with Proposed 765 kV Substations}
\end{figure}

%% ============================================================
\section{Generation (Operable EIA 860)}
%% ============================================================

\begin{table}[H]
\centering
\caption{Operable Generation Summary - Added in synthetic substations}
\begin{tabular}{lr}
\toprule
\textbf{Metric} & \textbf{Value} \\
\midrule
Total Gen Units      & 1,551 \\
Total Gen Fragments  & 1,136 \\
Total Capacity       & 147,977 MW \\
\bottomrule
\end{tabular}
\end{table}

\noindent\textit{Reference: EIA Form 860 (2024) --- operable generators in ERCOT region}

%% ============================================================
\section{Generation (GIS Planned --- IA Signed Approved)}
%% ============================================================

\begin{table}[H]
\centering
\caption{Planned Generation by Technology - Added in synthetic substations}
\begin{tabular}{lrr}
\toprule
\textbf{Technology} & \textbf{Units} & \textbf{MW} \\
\midrule
Solar Photovoltaic   & 252 & 60,152 \\
Battery Storage      & 209 & 36,598 \\
Onshore Wind         & 78  & 17,174 \\
Natural Gas (Steam)  & 16  & 5,599 \\
Other                & 10  & 925 \\
\midrule
\textbf{Total Planned} & \textbf{565} & \textbf{120,448} \\
\bottomrule
\end{tabular}
\end{table}

\noindent\textit{Reference: ERCOT GIS Interconnection Queue Report (January 2026 vintage) --- filtered to IA Signed status}

%% ============================================================
\section{Combined Generation by Fuel Type (2031 Case)}
%% ============================================================

\begin{table}[H]
\centering
\caption{Combined Generation by Fuel Type (EIA Operable + GIS Planned)}
\begin{tabular}{lrrrr}
\toprule
\textbf{Fuel} & \textbf{Units} & \textbf{Synthetic (MW)} & \textbf{Reference (MW)} & \textbf{Diff (\%)} \\
\midrule
Biomass              & 1   & 100    & 131    & $-$23.7 \\
Coal                 & 26  & 14,598 & 11,596 & +25.9 \\
Diesel               & 296 & 629    & 333    & +88.9 \\
Gas, Combined Cycle  & 266 & 39,649 & 33,341 & +18.9 \\
Gas Turbine          & 396 & 11,692 & 9,723  & +20.3 \\
Gas, Steam Turbine   & 72  & 17,586 & 10,383 & +69.4 \\
Hydro                & 25  & 472    & 569    & $-$17.0 \\
Nuclear              & 5   & 4,994  & 4,973  & +0.4 \\
Solar                & 380 & 77,899 & 94,375 & $-$17.5 \\
Battery Storage      & 318 & 42,116 & 61,258 & $-$31.2 \\
Wind                 & 257 & 50,328 & 50,006 & +0.6 \\
Other                & 17  & 357    & 0      & N/A \\
\midrule
\textbf{Total}       & \textbf{2,059} & \textbf{260,420} & \textbf{276,688} & \textbf{$-$5.9} \\
\bottomrule
\end{tabular}
\end{table}

\noindent\textit{Synthetic: EIA operable (2024) + GIS planned (IA Signed, Jan 2026). Reference: 2031 ERCOT case target.}

%% ============================================================
\section{Load by Voltage Level}
%% ============================================================

\begin{table}[H]
\centering
\caption{Load by Voltage Level}
\begin{tabular}{rrr}
\toprule
\textbf{kV} & \textbf{Loads} & \textbf{MW} \\
\midrule
69  & 1,605 & 34,683 \\
138 & 2,626 & 78,539 \\
345 & 48    & 32,351 \\
\bottomrule
\end{tabular}
\end{table}

%% ============================================================
\section{Large Load Allocation (2031 CDR Target)}
%% ============================================================

\begin{table}[H]
\centering
\caption{Large Load Allocation}
\begin{tabular}{lr}
\toprule
\textbf{Metric} & \textbf{Value} \\
\midrule
Total Large Load Target  & $\sim$52,350 MW \\
Large-Load Substations   & $\sim$250 (tunable, each $\geq$ 75 MW) \\
Avg Load Size            & $\sim$208 MW (benchmark, not scaled) \\
\bottomrule
\end{tabular}
\end{table}

\noindent\textit{Number of large loads is controlled by \texttt{N\_LARGE\_LOADS} (default $\sim$250). Each load keeps its benchmark MW size from per-zone lognormal sampling fitted to the ERCOT 2026 reference case (Thomas Cluster Process). No MW scaling is applied.}

\begin{figure}[H]
\centering
\includegraphics[width=0.85\textwidth]{figures/tsp-large-load-breakdown.png}
\caption{2031 ERCOT CDR Large Load Forecast}
\end{figure}

\subsection*{Large Load Design Assumptions}

Based on publicly available ERCOT interconnection queue data ($\sim$238,000 MW of requests):

\begin{table}[H]
\centering
\caption{Large Load Types (ERCOT Queue Proportions)}
\small
\begin{tabular}{clrl}
\toprule
\textbf{\#} & \textbf{Type} & \textbf{\% of MW} & \textbf{Dispatch Assumptions} \\
\midrule
1 & Data Centers          & 77.6\% & Firm, non-interruptible baseload; constant-power \\
2 & None / No type        & 11.1\% & Expected to be data centers; apply DC assumptions \\
3 & Crypto Mining         & 6.8\%  & Price-elastic CLR with bid curves; shed within minutes \\
4 & Industrial            & 2.7\%  & Mostly firm; 10--20\% curtailable at high prices \\
5 & DC / Crypto mixed     & 1.1\%  & Apply crypto dispatch assumptions \\
6 & Hydrogen Electrolysis & 0.6\%  & Flexible with ramp constraints \\
7 & Unknown               & 0.1\%  & Apply DC and crypto assumptions proportionally \\
\bottomrule
\end{tabular}
\end{table}

\noindent For loads under 75 MW, the standard census/population distribution method is used.

\subsection*{Current Large Load Assignment}

\begin{table}[H]
\centering
\caption{Current Synthetic Case Large Load Distribution}
\begin{tabular}{lrrr}
\toprule
\textbf{Type} & \textbf{Subs} & \textbf{MW} & \textbf{\% of Total} \\
\midrule
DataCenter              & 88  & 12,563 & 24.0\% \\
LargeLoad (unassigned)  & 158 & 39,787 & 76.0\% \\
\midrule
\textbf{Total}          & \textbf{246} & \textbf{52,350} & \textbf{100\%} \\
\bottomrule
\end{tabular}
\end{table}

\noindent\textit{DataCenter labeling is based on proximity to real TX data center coordinates from the IM3 atlas (within 50 km). All 92 IM3 TX data centers are located east of $-98.8^{\circ}$ longitude --- no data centers exist in the West or Far West zones.}

\begin{figure}[H]
\centering
\includegraphics[width=0.85\textwidth]{figures/765kV_design_overview_panel3.png}
\caption{Large Load Centers by Type and 765 kV Substations}
\end{figure}

\noindent\textit{\textbf{Note:} Using the IM3 Open Source Data Center Atlas (50 km proximity), only $\sim$22\% of large load MW is currently labeled as DataCenter. According to ERCOT interconnection queue data, data centers account for $\sim$77\% of large load demand. Closing this gap would require additional geographic data sources for data center locations beyond what IM3 currently provides.}

%% ============================================================
\section{765 kV Transmission Corridors}
%% ============================================================

Based on ERCOT's long-range 765 kV transmission plan. Existing 345 kV substations nearest to each corridor waypoint are upgraded. Waypoints are offset $\sim$20--50 km from ERCOT's published locations to maintain the synthetic nature of the grid.

\begin{table}[H]
\centering
\caption{765 kV Corridor Summary}
\begin{tabular}{llr}
\toprule
\textbf{Route} & \textbf{Segment} & \textbf{Substations} \\
\midrule
Permian Basin & West (NM border)    & 3 \\
Permian Basin & Middle (central TX) & 5 \\
Eastern       & DFW to Gulf Coast   & 4 \\
\midrule
\textbf{Total} & & \textbf{12} \\
\bottomrule
\end{tabular}
\end{table}

\noindent\textit{Reference: ERCOT Permian Basin Transmission Plan (765 kV corridors)}

\begin{figure}[H]
\centering
\includegraphics[width=\textwidth]{figures/Permian_basin_765KV_route.png}
\caption{ERCOT 765 kV Reference Plan}
\end{figure}

\begin{figure}[H]
\centering
\includegraphics[width=0.85\textwidth]{figures/765kV_design_overview_panel2.png}
\caption{Proposed 765 kV Substation Locations}
\end{figure}

\textbf{Permian Basin Route:} Far west Texas (NM border) through the Permian Basin to central Texas and the DFW area. Serves renewable generation export and west Texas load growth.

\textbf{Eastern Route:} DFW hub east toward Arkansas/Louisiana border and south along the Gulf Coast. Serves data centers, industrial loads, and coastal population centers.

\end{document}
